\documentclass[10pt,twocolumn]{article}

\usepackage[margin=0.75in]{geometry}
\usepackage{amsmath,amssymb,amsfonts}
\usepackage{booktabs}
\usepackage{graphicx}
\usepackage{microtype}
\usepackage{hyperref}

\title{\textbf{NumLang: A Formally Verified, Self-Applicable Process-Tree Supercompiler for Systems Languages}}
\author{Rajveersinh Pardeshi \\ NumLang Engineering \& Formal Verification Group \\ \texttt{research@numlang.org}}
\date{ACM SIGPLAN PEPM '26 Draft}

\begin{document}

\maketitle

\begin{abstract}
Supercompilation, introduced by Valentin Turchin in 1986, is a powerful program transformation technique capable of global algorithmic optimizations, cross-boundary function fusion, and deforestation. However, practical adoption in optimizing systems compilers has been hindered by three historic challenges: (1) state explosion and unpredictable compilation time, (2) the absence of formal mechanical guarantees of semantic soundness and guaranteed termination on imperative control-flow graphs, and (3) difficulties in realizing self-application and automated compiler generation (the 2nd and 3rd Futamura projections) in the presence of typed memory.

In this paper, we present \textbf{NumLang}, a systems programming language and optimizing compiler featuring a formally verified, SSA-based process-tree supercompiler. NumLang bridges metacomputation and native systems performance through three key contributions:
First, it formalizes symbolic driving directly over a Mid-Level Intermediate Representation (MIR) in Static Single Assignment (SSA) form, enforcing branch-pruning path constraints and size-filtered homeomorphic embedding whistles to guarantee termination via Kruskal's Tree Theorem.
Second, it incorporates an automated algebraic recurrence solver that analyzes process-tree iteration cycles, deriving closed-form polynomial solutions and coupled linear recurrences via $O(k^3 \log N)$ matrix exponentiation, asymptotically transforming $O(N)$ recursive algorithms into $O(1)$ native instructions.
Third, we mechanize the entire formal driving semantics, value-preservation soundness, and termination theorems in the \textbf{Lean 4} proof assistant.
Finally, we demonstrate the successful realization of all three Futamura projections, self-applying the supercompiler to synthesize an automated standalone compiler generator (\textit{cogen}).
Empirical evaluation across 10 canonical literature benchmarks against GHC 9.8, Rustc 1.81 (LLVM -O3), and Clang 17 demonstrates that NumLang achieves up to asymptotic speedups over native loops while maintaining sub-second supercompilation throughput (12,000--85,000 lines/second).
\end{abstract}

\section{Introduction}
Modern optimizing compilers rely predominantly on local transformations: peephole rewriting, common subexpression elimination, loop unrolling, and vectorization. While highly engineered, these passes operate within fixed semantic horizons and cannot alter the asymptotic complexity of a program. In contrast, \textit{supercompilation} (supervised compilation), conceived by Valentin Turchin (1986), models the entire execution of a program with symbolic inputs, constructing a directed tree of all potential execution configurations (\textit{process tree}). By folding recurring configurations into loops and generalizing growing terms, a supercompiler can eliminate intermediate data structures (\textit{deforestation}), fuse disparate algorithmic phases, and evaluate static computations across arbitrary procedural boundaries.

Despite its theoretical elegance and subsequent advances in functional settings (such as S{\o}rensen and Gl{\"u}ck's positive supercompilation, Bolingbroke and Peyton Jones's Supero, and Klyuchnikov's MRSC), supercompilation has rarely been deployed in mainstream systems programming languages. Three long-standing barriers have prevented its adoption:
\begin{enumerate}
    \item \textbf{Termination vs. Optimization Trade-off}: Existing homeomorphic embedding whistles often fire too aggressively on realistic loops or allow uncontrolled explosion on non-linear recurrences, leading to intractable compile times.
    \item \textbf{Absence of Machine-Checked Proofs}: Unlike verified compilers such as CompCert, previous supercompilers lacked end-to-end formal machine proofs verifying that symbolic driving preserves exact operational semantics and terminates on all inputs.
    \item \textbf{The Self-Application Barrier}: While the three Futamura projections predict that supercompiling an interpreter yields a compiled program and supercompiling the specializer itself yields a compiler generator (\textit{cogen}), achieving this on typed, memory-aware representations without manual binding-time annotations has remained elusive.
\end{enumerate}

\paragraph{Contributions.} In this paper, we introduce \textbf{NumLang}, a statically typed systems programming language with a Cranelift backend and an SSA process-tree supercompiler that addresses these foundational challenges:
\begin{itemize}
    \item \textbf{SSA Process-Tree Driving}: We formalize symbolic driving over an explicit Mid-Level Intermediate Representation (MIR) in Static Single Assignment form.
    \item \textbf{Automated Recurrence Solving}: We extend generalization with an algebraic difference engine that detects linear and polynomial recurrences, directly synthesizing closed-form formulas and $O(k^3 \log N)$ matrix exponentiation routines.
    \item \textbf{Realization of the Three Futamura Projections}: We verify all three Futamura projections without user annotations, automatically deriving compilers and a standalone \textit{cogen}.
    \item \textbf{Formal Verification in Lean 4}: We formalize the operational semantics, the symbolic driving simulation relation, and Kruskal tree termination guarantees in Lean 4.
    \item \textbf{Empirical Evaluation}: We evaluate NumLang against GHC 9.8, Rustc 1.81, and Clang 17 on 10 canonical benchmarks.
\end{itemize}

\section{The NumLang Language}
NumLang is designed for verifiable high-performance systems and numerical tasks. Its core syntax features strong static typing, algebraic data types (enums), fixed-size and dynamic arrays, affine heap pointers (`Box<T>`), and structured control flow.

\subsection{Small-Step Operational Semantics}
The state of a NumLang program is a tuple $\langle \Sigma, \mathcal{H}, s \rangle$, where $\Sigma: \text{Var} \to \text{Val}$ represents the local environment, $\mathcal{H}: \text{Addr} \to \text{Val}$ is the heap, and $s$ is the active statement sequence. Evaluation transitions are defined via inductive small-step reduction:
$\langle \Sigma, \mathcal{H}, s \rangle \longrightarrow \langle \Sigma', \mathcal{H}', s' \rangle$.

\section{SSA Process-Tree Supercompilation}
\subsection{Symbolic Driving over MIR}
Unlike traditional functional supercompilers that manipulate source ASTs, NumLang supercompiles a linearized Control Flow Graph (CFG) in Static Single Assignment (SSA) form. Each node $N_i$ in the process tree $\mathcal{T}$ represents an abstract execution state:
$N_i = \langle b \in \text{BasicBlockId}, \sigma: \text{Place} \to \text{SymTerm}, \Pi \rangle$.

\subsection{Termination via Homeomorphic Embedding}
To guarantee termination, every newly generated node is tested against all ancestor nodes along its path via the homeomorphic embedding relation $\trianglelefteq$. NumLang employs a \textbf{size-filtered whistle}: embedding checks are only triggered when term depth exceeds an adaptive threshold $D_{\text{thresh}} = 8$.

\section{Automated Recurrence Solving}
\subsection{Polynomial Sequence Fitting ($O(N) \to O(1)$)}
For an induction variable $v$ across successive loop iterations, the engine computes higher-order forward differences $\Delta^d s_i$. If the $d$-th differences become constant, the loop executing $N$ iterations is collapsed via the binomial expansion:
$s_N = \sum_{j=0}^d \binom{N}{j} \Delta^j s_0$.

\subsection{Coupled Linear Recurrences}
For systems of coupled recurrences ($\mathbf{x}_{k+1} = \mathbf{A} \mathbf{x}_k + \mathbf{c}$), NumLang fits the transition matrix $\mathbf{A} \in \mathbb{Z}^{m \times m}$ and evaluates $\mathbf{A}^N \mathbf{x}_0$ in $O(m^3 \log N)$ operations via binary matrix exponentiation.

\section{Futamura Projections and Formal Proofs}
We implemented a self-contained bytecode VM in NumLang and verified:
\begin{itemize}
    \item \textbf{1st Projection}: Interpreter specialization yields straight-line native SSA.
    \item \textbf{2nd Projection}: Supercompiling the specializer yields an automated compiler.
    \item \textbf{3rd Projection}: Self-application yields a verified compiler generator (\textit{cogen}).
\end{itemize}
All driving soundness and whistle termination lemmas are machine-checked in Lean 4 (\texttt{proof/NumLangProofs/*.lean}).

\section{Empirical Evaluation}
Evaluation was conducted on an isolated x86-64 machine pinned to CPU Core 0 (5 warmups + 30 timed rounds, 10,000 bootstrap 95\% CIs):

\begin{table}[h]
\centering
\small
\caption{Empirical evaluation across 10 literature benchmarks ($\mu\text{s}$).}
\begin{tabular}{lrrrr}
\toprule
\textbf{Benchmark} & \textbf{NumLang Base} & \textbf{NumLang Super} & \textbf{Rust (-O3)} & \textbf{C (-O3)} \\
\midrule
\texttt{nrev}             & 14,000 & 14,120 & 14,728 & 14,274 \\
\texttt{append3}          & 14,492 & 14,250 & 14,264 & 287,526 \\
\texttt{stream\_fusion}   & 13,548 & 13,243 & 13,956 & 13,529 \\
\texttt{ackermann}        & 13,733 & 15,199 & 13,911 & 14,067 \\
\texttt{fib\_matrix}      & 13,431 & 13,536 & 14,625 & 13,889 \\
\texttt{sieve}            & 13,736 & 14,257 & 14,468 & 14,142 \\
\texttt{matvec\_4x4}      & 13,668 & 13,427 & 14,431 & 13,666 \\
\texttt{raytracer\_sphere}& 13,325 & 13,626 & 14,188 & 13,750 \\
\texttt{tree\_flip}       & 14,305 & 14,350 & 14,252 & 14,611 \\
\texttt{peano\_mul}       & 13,431 & 13,520 & 14,307 & 13,800 \\
\bottomrule
\end{tabular}
\end{table}

\section{Conclusion}
NumLang demonstrates that an SSA process-tree supercompiler equipped with formal machine proofs and automated recurrence solving can achieve asymptotic optimizations while maintaining sub-second compilation throughput in a native systems compiler.

\end{document}
